% from ega0-19-fr.tex
\oldpage[0\textsubscript{IV-1}]{69}

\begin{env}[19.0.1]
\label{0.19.0.1-fr}
Au chapitre IV, nous introduirons et étudierons entre autres une classe
importante de morphismes de préschémas, les morphismes \emph{lisses}\footnote{Aussi
appelés morphismes \emph{simples} dans certains travaux récents (en s'inspirant
de la terminologie classique de «~point simple~»)~; cette terminologie prête
cependant à confusion, notamment en théorie des groupes algébriques.}. Une de
leurs propriétés fondamentales (et qui, jointe à une condition de finitude, les
caractérise) est une propriété de \emph{relèvement des morphismes}~: si
$f:X\to Y$ est un morphisme lisse, $g:Y'\to Y$ un morphisme, alors pour tout
morphisme $h:Y''\to Y'$ faisant de $Y''$ un préschéma «~peu différent~» de
$Y'$, \emph{tout $Y$-morphisme $Y''\to X$ se factorise} en
$Y''\xrightarrow{h}Y'\to X$. De façon précise, bornons-nous au cas où
$Y=\operatorname{Spec}(A)$, $X=\operatorname{Spec}(B)$ sont \emph{affines}~;
$B$ est alors appelée une $A$-algèbre \emph{formellement lisse} si, pour toute
$A$-algèbre $C$ et tout idéal \emph{nilpotent} $\mathfrak J$ de $C$, tout
$A$-homomorphisme $B\to C/\mathfrak J$ \emph{se relève} en
$B\to C\to C/\mathfrak J$. En d'autres termes, l'application
\[
  \operatorname{Hom}_A(B,C)\longrightarrow
  \operatorname{Hom}_A(B,C/\mathfrak J)
\]
est \emph{surjective}. Dans beaucoup d'applications, $X$ apparaîtra comme objet
représentant un \emph{foncteur contravariant représentable} $Y'\mapsto F(Y')$
de la catégorie des $Y$-préschémas dans celle des ensembles, de sorte
($\mathbf{0}_{\mathrm{III}}$, 8.1.8) que l'on aura
$F(Y')=\operatorname{Hom}_Y(Y',X)$. Dans le cas affine, si l'on pose
$F^0(C)=F(\operatorname{Spec}(C))$, la vérification du fait que, sous les
conditions ci-dessus, $F^0(C)\to F^0(C/\mathfrak J)$ est \emph{surjectif}
(qui peut se faire même \emph{sans savoir} que $F$ est représentable) établira
que $f$ est lisse, pourvu que la condition additionnelle de finitude soit
remplie.
\end{env}

\begin{env}[19.0.2]
\label{0.19.0.2-fr}
Pour les \emph{algèbres topologiques} sur un anneau topologique $A$, il y a une
notion analogue d'algèbre \emph{formellement lisse} que nous ne préciserons pas
ici (cf.~19.3.1). L'étude de ces notions est faite d'abord d'un point de vue
élémentaire au §~19, puis, à l'aide des propriétés des différentielles qui
seront développées aux §§~20 et 21, des théorèmes plus délicats seront prouvés
au §~22. Résumons ici les principaux résultats sur les algèbres formellement
lisses~:

\medskip\noindent\textbf{I.} --- Soient $A$, $B$ deux anneaux locaux
noethériens, $\varphi:A\to B$ un homomorphisme local, $k$ le corps résiduel de
$A$, et soit $B_0=B\otimes_A k$~; on munit $A$, $B$ et $B_0$ de leurs
topologies préadiques et $k$ de la topologie discrète. Alors, pour que $B$ soit
une $A$-algèbre formellement lisse, il faut et il suffit que $B$ soit un
$A$-module \emph{plat} et que $B_0$ soit une $k$-algèbre \emph{formellement
lisse} (19.7.1). Ce théorème ramène donc la lissité formelle, pour les anneaux
locaux noethériens, à la même question pour les anneaux locaux noethériens qui
sont des algèbres sur un \emph{corps}.

\medskip\noindent\textbf{II.} --- Soient $k$ un corps, $A$ un anneau local
noethérien qui est une $k$-algèbre. Pour que $A$ soit formellement lisse, il
faut et il suffit que $A$ soit \emph{géométriquement régulier} sur $k$,
c'est-à-dire que pour toute extension \emph{finie} $k'$ de $k$, l'anneau
semi-local $A\otimes_k k'$ soit

